%MIT OpenCourseWare: https://ocw.mit.edu
%RES.18-012 Algebra II Student Notes, Spring 2022
%License: Creative Commons BY-NC-SA 
%For information about citing these materials or our Terms of Use, visit: https://ocw.mit.edu/terms.

% \rhead{\rightmark}
\lhead{Lecture 8: Rings}

\section{Rings}

We'll now discuss the second main topic of this course: rings.

\begin{qq}
Groups have only one operation, but lots of familiar sets ($\ZZ, \RR, \CC, \QQ$) have two operations: addition and multiplication. How can we generalize the idea of sets with two operations, rather than one?
\end{qq}

\subsection{What is a Ring?}

Informally, a ring is a set of elements which can be added and multiplied, so that the natural properties we would expect of addition and multiplication all hold. 

\begin{example}
The familiar sets $\ZZ, \QQ, \RR,$ and $\CC$ are rings. 
\end{example}

\begin{example}
There are various rings between $\ZZ$ and $\QQ$: for example, \[\ZZ[1/2] \coloneqq \left\{a/2^k \mid a, k \in \ZZ\right\}\] is a ring. Similarly, the \textbf{Gaussian integers} \[\ZZ[i] = \{a + bi \mid a, b \in \ZZ\}\] and \[\QQ[\sqrt{2}] = \{a + b\sqrt{2} \mid a, b \in \QQ\}\] are rings. So is \[\ZZ[\zeta] = \left\{a_0 + a_1\zeta + \cdots + a_{n - 1}\zeta^{n - 1} \mid a_i \in \ZZ\right\},\] where $\zeta$ is some $n$th root of unity. 
\end{example}

We can now state the formal definition of a ring:

\begin{definition}
A \textbf{ring} $R$ is a set with two binary operations $R \times R \to R$, written as $+$ and $\cdot$ (and called addition and multiplication), which satisfy the following axioms: 
    \begin{enumerate}
        \item $(R, +)$ is an abelian group: addition is associative and commutative, there is an additive identity $0_R$ such that $0_R + a = a$ for all $a \in R$, and every element has an additive inverse. 
        \item Multiplication is also associative and commutative, and there is a multiplicative identity $1_R$. (In other words, under multiplication, $R$ is a \emph{semigroup} --- a group without the condition that every element has an inverse.)
        \item Addition and multiplication satisfy \emph{distributivity}: for all $a, b, c \in R$, we have \[a(b + c) = ab + ac.\] 
    \end{enumerate}
\end{definition}
    
% A \textbf{ring} is a set with two binary operations $R \by R \rto R$, written as $+: (a, b) \mto a+b$ and $\cdot : (a, b) \mto a\cdot b$ or $ab.$ The operations must satisfy:
% \begin{enumerate}[label = (\Roman*)]
%     \item Addition is commutative\todo{add these operations}
% \end{enumerate}

In this class, we'll use ``ring'' in this sense; but usually a ring satisfying this definition is called a \emph{commutative unital ring}. In defining a general ring, one can drop the requirement of commutativity of multiplication, or the existence of $1_R$. If every condition holds except for $ab = ba$ --- and we add distributivity in the other direction as well, meaning $(b + c)a = ba + ca$ --- then $R$ is called a \textbf{noncommutative ring}. 


\begin{example}[Matrices]
The ring of matrices $\matnn(\CC)$ is a noncommutative ring, since matrix multiplication is noncommutative (and all the other axioms are satisfied). Similarly, $\edo(V)$ for a vector space $V$ is a noncommutative ring.
\end{example}

\begin{example}[Group Ring]
If $G$ is a group, then take the vector space with basis vectors $v_g$ for $g \in G$ (note that this is the vector space acted on by the regular representation). Clearly, addition is already defined; and multiplication can be defined in the natural way, as \[v_gv_h = v_{gh}.\] Then this gives the \textbf{group ring}, which is noncommutative if $G$ is nonabelian. 
\end{example}

From now on, all rings will be commutative as in the definition, unless stated otherwise. 

\subsection{Zero and Inverses}

The following proposition confirms a property that we would like rings to have.

\begin{proposition}
In any ring $R$, $0_R \cdot a = 0_R$ for all $a \in R$.
\end{proposition}

\begin{proof}
Pick some $x \in R$. Then, since $0_R + x = x$, we have \[xa = (0_R + x)a = 0_Ra + xa.\] We can cancel out $xa$ (since $R$ is an abelian group under addition), so $0_R = 0_Ra.$
\end{proof}

\begin{corollary}
The additive identity $0_R$ cannot have a multiplicative inverse unless $0_R = 1_R$. (In other words, division by $0$ is only possible when $0 = 1$.)
\end{corollary}

The axioms do not require that $0_R \neq 1_R$. But if $0_R = 1_R$, then for any $x \in R$, \[x = x\cdot 1_R = x\cdot 0_R = 0_R.\] So $R$ must be a one-element ring; there's only one binary operation on a set with one element, which does satisfy the axioms. This ring is called the \textbf{zero ring}; it is a legitimate but trivial example. In all other cases, $0_R \neq 1_R$, as we would expect. 

\begin{definition}
A (nonzero) ring where every nonzero element has a multiplicative inverse is called a \textbf{field}.
\end{definition}

Note that by definition, the zero ring is not a field. 

\begin{example}
The familiar sets $\QQ$, $\RR$, and $\CC$ are fields. 
\end{example}

\begin{example}
The integers $\ZZ$ do not form a field, since most numbers do not have inverses. For the same reason, the Gaussian integers $\ZZ[i]$ are not a field either. For example, $2$ is not invertible in either ring.
\end{example}

\begin{example}
The integers modulo $n$, denoted $\ZZ/n\ZZ$, form a field if and only if $n$ is prime (since in general, $a$ is invertible mod $n$ if and only if $a$ and $n$ are coprime). 
\end{example}

More examples of rings (which are not fields) come from functions:

\begin{example}
    The set $\mathbb{C}[x]$, consisting of polynomials in one variable with complex coefficients, is a ring. The set $C^\infty(\mathbb{R})$, consisting of real-valued functions which are continuous and infinitely differentiable, is also a ring. 
\end{example}

In some sense, you can think of fields as a generalization of numbers, and more general rings as generalizations of functions. 

\subsection{Homomorphisms}

When studying groups, one of the most powerful tools comes from thinking about mappings between groups. In particular, homomorphisms, which are functions that respect the group operation, and their kernels/images provide lots and lots of information about groups. 

\begin{qq}
How can we formalize the idea of "functions between rings that behave nicely with respect to addition and multiplication"?
\end{qq}

Similarly to the path we took when studying groups, we can now define a ring homomorphism.

 \begin{definition}
 A ring \textbf{homomorphism} from $R$ to $S$ is a map $\varphi: R \to S$ such that:
 \begin{enumerate}
     \item $\varphi(a+b) = \varphi(a) + \varphi(b)$ for all $a, b \in R$;
     \item $\varphi(ab) = \varphi(a)\varphi(b)$ for all $a, b \in R$;
     \item $\varphi(1_R) = 1_S.$
 \end{enumerate}
 \end{definition}
 
 \begin{example}
The mapping $\ZZ \to \ZZ/n\ZZ$ given by $a \mapsto \overline{a}$ (where $\overline{a}$ denotes $a$ mod $n$) is a ring homomorphism.
\end{example}

Why is it not necessary to require that $\varphi(0_R) = 0_S$? In fact, it is implied by the other axioms, because additive inverses exist! 

\begin{proposition}
For a ring homomorphism $\varphi: R \to S$, it must be the case that $\varphi(0_R) = 0_S$. 
\end{proposition}

\begin{proof}
Using the first property, $\varphi(a) = \varphi(a + 0_R) = \varphi(a) + \varphi(0_R)$, and so adding $-\varphi(a)$ to both sides gives $0_S = \varphi(0_R)$.
\end{proof}

On the other hand, because there are not necessarily multiplicative inverses, the property $\varphi(1_R) = 1_S$ must be explicitly written. In fact, there are examples of maps compatible with $+$ and $\cdot$ (meaning they satisfy the first two properties) that have $\varphi(1_R) \neq 1_S$.

\begin{example}
Let $R$ be the zero ring, and $S$ be any nonzero ring (for example, $\ZZ$). Then take the map $0_R \mapsto 0_S$. This is compatible with the additive and multiplicative structure of the two rings; but since $1_S \neq 0_S$, it is not a ring homomorphism. (There exist less trivial examples as well.)
\end{example}

If $\varphi : R \to S$ is one-to-one and onto (that is, $\varphi$ is a bijection), then it is possible to check (as in the case of group homomorphisms) that the inverse bijection is also a homomorphism.

\begin{definition}
A bijective homomorphism is called an \textbf{isomorphism}.
\end{definition}

If $\varphi$ is one-to-one but \emph{not} onto, then $\varphi$ is an isomorphism from $R$ to its image, so it identifies $R$ with a \textbf{subring} of $S$:

\begin{definition}\todo{explain this defintion/flesh out}
A \textbf{subring} $S$ of $R$ is a subset which is closed under addition, taking additive inverses, and multiplication, and contains $1_R$. 
\end{definition}

Meanwhile, if $\varphi$ is \emph{not} one-to-one, then we can think about its kernel \[\ker(\varphi) = \{a \in R: \varphi(a) = 0\}.\] (This is the same definition as we saw with groups.) If $\varphi$ is not one-to-one, then $\ker(\varphi)$ is nontrivial. 

\begin{qq}
What information can be gained from ring homomorphisms with nontrivial kernel?
\end{qq}

The kernel must be an additive subgroup of $S$. But it turns out that it must also be compatible with multiplication in some ways; this brings us to the concept of \emph{ideals}. 

\subsection{Ideals}

\vspace{0.75em}

\begin{definition}[Ideal]
An \textbf{ideal} of $R$ is a subset $I \subset R$ such that $I$ is an additive subgroup of $R$, and for any $x \in I$ and $a \in R$, we have $ax \in I$. 
\end{definition}

So an ideal isn't just closed under multiplication, it's in fact closed under multiplication by \emph{any} element in $R$. 

The kernel of a ring homomorphism is necessarily an ideal --- if $\varphi(x) = 0$, then for any $a \in R$, we have \[\varphi(ax) = \varphi(a)\varphi(x) = 0.\]  

\begin{example}
For any $n \in \ZZ$, the subset $n\ZZ$ (consisting of multiples of $n$) is an ideal. More generally, if $R$ is any ring and $a$ an element of $R$, then $aR = \{ax \mid x \in R\}$ is an ideal. 
\end{example}

In fact, this is an important example of an ideal, as we'll see later, so it has a name:

\begin{definition}[Principal Ideal]
For an element $a \in R$, the ideal $aR$ is called a \textbf{principal ideal}, and denoted $(a)$. 
\end{definition}

Any additive subgroup of $\ZZ$ is cyclic, i.e., of the form $n\ZZ$. This means every ideal of $\ZZ$ is principal (since every ideal is an additive subgroup). However, in a general ring, not every ideal will be principal. 

More generally, we can consider picking \emph{several} elements:
\begin{definition}
For elements $a_1, \ldots, a_n \in R$, the set of linear combinations \[\left\{\sum a_ix_i \mid x_i \in R\right\}\] is an ideal of $R$, denoted as $(a_1, \ldots, a_n)$; this is called the ideal \textbf{generated} by $a_1$, \ldots, $a_n$. 
\end{definition}

Note that $(a_1, \ldots, a_n)$ is the smallest ideal containing all of $a_1$, \ldots, $a_n$ (as it is an ideal, while all elements $\sum a_ix_i$ must be in the ideal by the axioms). 

% Consider picking several elements $a_1, \cdots, a_n \in R.$ The set of linear combinations is \[(a_1, \cdots, a_n) \coloneqq \{\sum a_ix_i: x_I \in R\},\] and is an ideal. It is called the ideal \textbf{generated} by $a_1, \cdots, a_n$, and is the smallest ideal containing each $a_i.$\footnote{All of these concepts are analogous to concepts studied in group theory (in 18.701).} 
Ideals in rings are in some sense analogous to normal subgroups in groups (in particular, both arise as the kernel of a homomorphism), and we can take quotients in a similar way as well:

\begin{proposition}[Quotient Ring]
Let $R$ be a ring and $I \subset R$ an ideal. Since $I$ is a normal subgroup of $R$ under addition (as both are abelian groups), we can construct the quotient $R/I$ of additive groups. Then $R/I$ is in fact a ring (with multiplication defined in the natural way --- the product of the cosets corresponding to $x$ and $y$ is the coset corresponding to $xy$), called the \textbf{quotient ring}. 
\end{proposition}

\begin{proof}
For each $x \in R$, we use $\overline{x}$ to denote the coset $x + I$. 

We first need to check that multiplication is well-defined, meaning that if $\overline{x}$ is represented by two elements $x_1$ and $x_2$, then using either representative to calculate $\overline{x}\cdot \overline{y}$ will give us the same result. But we have $x_1 - x_2 = a$ for some $a \in I$, which means \[x_1y - x_2y = ay \in I\] as well (since $I$ is closed under multiplication by any element $y \in R$), so $x_1y$ and $x_2y$ are also in the same coset. So the product of two cosets doesn't depend on our choice of representatives, which means multiplication really is well-defined. 

As in the case of groups, once we know that multiplication is well-defined, it is easy to check that all the ring axioms are satisfied by $R/I$. 
\end{proof}

Without going into detail, replacing ``group'' by ``ring'' and ``normal subgroup'' by ``ideal,'' the story for rings is extremely similar to that for groups. For example, if $\varphi$ is a homomorphism, then there is an isomorphism \[R/\ker(\varphi) \cong \im(\varphi).\] Additionally, for an ideal $I \subset R$, there is a bijection between ideals in $R/I$ and ideals in $R$ containing $I$ (this is essentially the Correspondence Theorem for rings). 

% If $\varphi$ is a homomorphism, then there is an isomorphism $R/\ker(\varphi) \cong \im(\varphi),$ and there is a correspondence between ideals $R/I$ and ideals in $R$ containing $I.$\footnote{This is the Correspondence Theorem for rings.\todo{check that this is true}}

\begin{example}
For $R = \ZZ$ and $I = n\ZZ,$ we have $R/I = \ZZ/n\ZZ$ (where $\ZZ/n\ZZ$ denotes the integers mod $n$) as rings, not just groups. This is essentially the fact that multiplication of residues mod $n$ is well-defined, not just addition. (This is where the notation $\ZZ/n\ZZ$ comes from --- to obtain the integers mod $n$, we're quotienting out the ring of integers by the ideal $n\ZZ$.) 
\end{example}